\documentclass[11pt, UTF8]{ctexart}
\usepackage{amsmath, amsfonts, amssymb}
\usepackage{extarrows} 
\usepackage{geometry} 
\usepackage{xcolor} 
\usepackage{enumitem} 
\usepackage[most]{tcolorbox} 
\usepackage{fancyhdr} 
\usepackage{diagbox} 
\usepackage{tikz}    
\setlength{\headheight}{13.6pt}
\geometry{a4paper, margin=0.8in} 
\pagestyle{fancy}
\fancyhf{}
\lhead{\color{fblue} \hwdate 作业} 
\rhead{\color{fblue} \today}
\cfoot{\thepage}

\definecolor{fblue}{RGB}{0, 74, 142}
\definecolor{fred}{RGB}{204, 26, 26}

\newtcolorbox{problem}[2]{
    enhanced,
    colback=white,
    colframe=fblue!80,
    arc=2mm,
    boxrule=0.8pt,
    before upper={\textbf{\color{fblue} 问题 #1 (#2)}\hspace{0.5em}},
    before skip=1.5em,
    after skip=1.5em,
    left=3mm, right=3mm, top=2mm, bottom=2mm
}

\newcommand{\hwdate}{5月14日} 
\newcommand{\Prob}{P}
\newcommand{\comp}[1]{\overline{#1}\mskip 3mu}
\newcommand{\ans}[1]{\mathbf{\color{fred}#1}}

\begin{document}

\begin{center}
    {\LARGE \bfseries \color{fblue} \hwdate 作业}
\end{center}

% ---------------- 题目 1 ----------------
\begin{problem}{1}{4. 28}
设二维随机变量 $(X,Y)$ 的概率密度为
\[ f(x,y) = \begin{cases} \frac{1}{\pi}, & x^2 + y^2 \leqslant 1, \\ 0, & \text{其他.} \end{cases} \]
试验证 $X$ 和 $Y$ 是不相关的，但 $X$ 和 $Y$ 不是相互独立的.
\end{problem}

\paragraph{\color{fblue}证} 
\begin{align*}
E(X) &= \int_{-\infty}^{\infty}\int_{-\infty}^{\infty} xf(x,y)\mathrm{d}x\mathrm{d}y = \iint\limits_{x^2+y^2 \leqslant 1} \frac{x}{\pi}\mathrm{d}x\mathrm{d}y \\
&= \frac{1}{\pi}\int_{-1}^{1} \mathrm{d}y\int_{-\sqrt{1-y^2}}^{\sqrt{1-y^2}} x\mathrm{d}x = 0.
\end{align*}
同样
\begin{align*}
E(Y) &= \int_{-\infty}^{\infty}\int_{-\infty}^{\infty} yf(x,y)\mathrm{d}x\mathrm{d}y = \iint\limits_{x^2+y^2 \leqslant 1} \frac{y}{\pi}\mathrm{d}x\mathrm{d}y = 0,
\end{align*}
而
\begin{align*}
E(XY) &= \int_{-\infty}^{\infty}\int_{-\infty}^{\infty} xyf(x,y)\mathrm{d}x\mathrm{d}y = \iint\limits_{x^2+y^2 \leqslant 1} \frac{xy}{\pi}\mathrm{d}x\mathrm{d}y \\
&= \frac{1}{\pi}\int_{-1}^{1} y\mathrm{d}y\int_{-\sqrt{1-y^2}}^{\sqrt{1-y^2}} x\mathrm{d}x = 0,
\end{align*}
从而 $\ans{E(XY) = E(X)E(Y),} $
这表明 $X,Y$ 是不相关的. 又
\[ 
f_X(x) = \int_{-\infty}^{\infty} f(x,y)\mathrm{d}y = \begin{cases} \int_{-\sqrt{1-x^2}}^{\sqrt{1-x^2}} \frac{1}{\pi}\mathrm{d}y = \frac{2}{\pi}\sqrt{1-x^2}, & -1 < x < 1, \\ 0, & \text{\textbf{其他.}} \end{cases} 
\]
同样
\[ 
f_Y(y) = \begin{cases} \frac{2}{\pi}\sqrt{1-y^2}, & -1 < y < 1, \\ 0, & \text{\textbf{其他.}} \end{cases} 
\]
显然 $\ans{f_X(x)f_Y(y) \neq f(x,y)}$, 故 $X,Y$ 不是相互独立的.

\vspace{1.5em}

% ---------------- 题目 2 ----------------
\begin{problem}{2}{4. 32}
设随机变量 $(X,Y)$ 具有概率密度
\[ f(x,y) = \begin{cases} \frac{1}{8}(x+y), & 0 \leqslant x \leqslant 2, 0 \leqslant y \leqslant 2, \\ 0, & \text{其他.} \end{cases} \]
求 $E(X),E(Y),\text{Cov}(X,Y),\rho_{XY},D(X+Y)$.
\end{problem}

\paragraph{\color{fblue}解}
注意到 $f(x,y)$ 在区域 $G: \{(x,y) \mid 0 < x < 2, 0 < y < 2\}$ 上不等于零，故有
\begin{align*}
E(X) &= \int_{-\infty}^{\infty}\int_{-\infty}^{\infty} xf(x,y)\mathrm{d}x\mathrm{d}y = \int_0^2 \mathrm{d}x\int_0^2 \frac{x}{8}(x+y)\mathrm{d}y \\
&= \int_0^2 \frac{x}{8}\left(xy + \frac{1}{2}y^2\right) \Bigg|_0^2 \mathrm{d}x = \int_0^2 \frac{x}{4}(x+1)\mathrm{d}x = \frac{7}{6}, \\
E(X^2) &= \int_{-\infty}^{\infty}\int_{-\infty}^{\infty} x^2f(x,y)\mathrm{d}x\mathrm{d}y = \int_0^2 \mathrm{d}x\int_0^2 \frac{x^2}{8}(x+y)\mathrm{d}y \\
&= \frac{1}{8}\int_0^2 x^2\left(xy + \frac{1}{2}y^2\right) \Bigg|_0^2 \mathrm{d}x = \frac{1}{4}\int_0^2 (x^3+x^2)\mathrm{d}x = \frac{5}{3}, \\
E(XY) &= \int_{-\infty}^{\infty}\int_{-\infty}^{\infty} xyf(x,y)\mathrm{d}x\mathrm{d}y = \int_0^2 \mathrm{d}x\int_0^2 \frac{xy}{8}(x+y)\mathrm{d}y \\
&= \frac{1}{4}\int_0^2 \left(x^2 + \frac{4x}{3}\right)\mathrm{d}x = \frac{4}{3}.
\end{align*}
由 $x,y$ 在 $f(x,y)$ 的表达式中的对称性(即在表达式 $f(x,y)$ 中将 $x$ 和 $y$ 互换，表达式不变)，得知
\[ E(Y) = E(X) = \ans{\frac{7}{6}}, \quad E(Y^2) = E(X^2) = \frac{5}{3}, \]
且有 
\[ D(Y) = D(X) = E(X^2) - [E(X)]^2 = \frac{5}{3} - \left(\frac{7}{6}\right)^2 = \frac{11}{36}. \]
而 
\begin{align*}
\text{Cov}(X,Y) &= E(XY) - E(X)E(Y) = \frac{4}{3} - \frac{49}{36} = \ans{-\frac{1}{36}}, \\
\rho_{XY} &= \frac{\text{Cov}(X,Y)}{\sqrt{D(X)}\sqrt{D(Y)}} = \ans{-\frac{1}{11}}, \\
D(X+Y) &= D(X) + D(Y) + 2\text{Cov}(X,Y) = \ans{\frac{5}{9}}.
\end{align*}

\vspace{1.5em}

% ---------------- 题目 3 ----------------
\begin{problem}{3}{4. 34}
(1) 设随机变量 $W = (aX+3Y)^2, E(X) = E(Y) = 0, D(X) = 4, D(Y) = 16, \rho_{XY} = -0.5$. 求常数 $a$ 使 $E(W)$ 为最小，并求 $E(W)$ 的最小值. \\
(2) 设随机变量 $(X,Y)$ 服从二维正态分布，且有 $D(X) = \sigma_X^2, D(Y) = \sigma_Y^2$. 证明当 $a^2 = \sigma_X^2/\sigma_Y^2$ 时，随机变量 $W = X-aY$ 与 $V = X+aY$ 相互独立.
\end{problem}

\paragraph{\color{fblue}解}
(1) 
\begin{align*} 
E(W) &= E[(aX+3Y)^2] = a^2E(X^2) + 6aE(XY) + 9E(Y^2), \\
E(X^2) &= D(X) + [E(X)]^2 = 4, \\
E(Y^2) &= D(Y) + [E(Y)]^2 = 16, \\
E(XY) &= \text{Cov}(X,Y) + E(X)E(Y) = \rho_{XY}\sqrt{D(X)D(Y)} = -4, 
\end{align*}
故
\[ E(W) = 4a^2 - 24a + 144 = 4(a-3)^2 + 108, \]
故当 $\ans{a = 3}$ 时 $E(W)$ 取最小值，$\min\{E(W)\} = \ans{108}$.

(2) 因为 $(X,Y)$ 是二维正态变量，而 $W$ 与 $V$ 分别是 $X,Y$ 的线性组合，故由 $n$ 维正态随机变量的性质 $3^{\circ}$ 知 $(W,V)$ 也是二维正态变量. 现在 $a^2 = \sigma_X^2/\sigma_Y^2$, 故有
\begin{align*}
\text{Cov}(W,V) &= \text{Cov}(X-aY, X+aY) \\
&= \text{Cov}(X,X) - a^2\text{Cov}(Y,Y) = \sigma_X^2 - a^2\sigma_Y^2 = \ans{0},
\end{align*}
即知 $W$ 与 $V$ 不相关. 又因 $(W,V)$ 是二维正态变量，故知 $W$ 与 $V$ 是相互独立的.

\end{document}